\label{chap.6}

A tester can only be evaluated on applications whose faults are known. This chapter
describes the four sets of test applications (Sections~\ref{sec:benchmark}--\ref{sec:external}),
how an application is admitted before any agent touches it
(Section~\ref{sec:admission}), how applications were split between development and
evaluation (Section~\ref{sec:split}), and the metrics and statistics
(Section~\ref{sec:stats}). Section~\ref{sec:threats} states the threats to validity.

\section{The six-application benchmark}
\label{sec:benchmark}

Six small web applications were built to one design contract for the v1 experiments: an
account hub, an admin records system, a booking desk, a media library, a shop catalogue and a
support portal. Each carries two seeded faults, 12 in total, in five classes: calculation
error, constraint bypass, persistence failure, state corruption and broken transition. Some
faults sit 9--16 contextual steps deep. Two applications (the shop catalogue and the
booking desk) were used during development; the other four stayed pristine until the
transfer experiment.

\section{The generated application family}

For Experiment 2, a seeded generator produced twelve variants of one loan-request workflow.
They differ in vocabulary, structure and decoys, and each was admitted by the same validator
as the benchmark (Section~\ref{sec:exp2}).

\section{The development suite}
\label{sec:devsuite}

Eight applications were built to be genuinely different from each other and from the
benchmark: different domains, architectures and stacks (Table~\ref{tab:devsuite}). They were
first built as zero-dependency applications. Then, at the project owner's request, they were
rebuilt as full-stack applications. Each now has its own database (PostgreSQL or MongoDB), a
sign-in page and a richer data model. The five rebuilt last were also enlarged to at least 12
views, 3 entities with full create, read, update and delete, 5 interaction patterns, and 2
new faults that are not about forms.

\begin{table}[H]
\centering
\small
\caption{The development suite (version 2).}
\label{tab:devsuite}
\begin{tabularx}{\textwidth}{@{}l X l l r r@{}}
\toprule
App & Domain and architecture & Server & Database & Faults & Crash \\
\midrule
tidewater & personal banking; single-page app over a JSON API & Python & PostgreSQL & 3 & -- \\
clinicflow & clinic front desk; server-rendered branching wizard & Python & PostgreSQL & 5 & -- \\
stackroom & stockroom inventory; data grid with server-side paging, sorting, bulk actions & Python & PostgreSQL & 5 & 2 \\
ledgerly & expense claims with multi-role approval; server-rendered forms & Node.js & PostgreSQL & 5 & -- \\
bagstreet & small shop; server-rendered pages and HTML forms & Python & MongoDB & 5 & -- \\
helpline & broadband care desk; pages with HTML fragment updates & Python & MongoDB & 7 & 2 \\
orbit & kanban boards; modal-heavy single-page app & Node.js & MongoDB & 5 & 2 \\
wayfarer & trip itinerary builder; client drafts and sealed server versions & Node.js & MongoDB & 6 & 2 \\
\midrule
& & & & 41 & 8 \\
\bottomrule
\end{tabularx}
\end{table}

\paragraph{Isolation between sessions.} Every signed-in browser context gets its own data
partition: a schema per session in PostgreSQL, a database per session in MongoDB. This
resets the application between episodes and keeps parallel runs independent. The database
drivers are vendored and pinned, so a build needs no network.

\paragraph{Crash-type faults.} Four applications also carry two crash-type faults each
(Section~\ref{sec:crash}), so that the tuned tester is exercised on the kind of failure the
external benchmark counts. By the project owner's decision only a subset of applications has
them. Orbit's were admitted (40 of 40 checks, including Docker); helpline, stackroom and wayfarer were
admitted with all checks passing (44, 40 and 41 checks, including Docker).

\section{The validation suite}
\label{sec:valsuite}

Eight further applications were planned by an independent agent working only from the
suite's design contract, without access to the tester's code or results, and built for the
held-out evaluation (Table~\ref{tab:valsuite}). Each has three seeded faults, 24 in total,
and each was admitted with 20 of 20 checks. They were evaluated in their first version
(without a database server or sign-in) and have not been changed since.

\begin{table}[H]
\centering
\small
\caption{The validation suite, used once, for the held-out evaluation.}
\label{tab:valsuite}
\begin{tabularx}{\textwidth}{@{}l X l r@{}}
\toprule
App & Domain & Stack & Faults \\
\midrule
dutyday & a food bank's volunteer rota & Python, SQLite & 3 \\
fosterly & an application to become a foster carer, saved section by section & Node.js, JSON files & 3 \\
larder & a personal recipe box with scaling and unit conversion & Node.js, JSON files & 3 \\
markbook & a teacher's markbook with a student view & Python, SQLite & 3 \\
quizcraft & a survey editor with undo, redo and a respondent preview & Python, SQLite & 3 \\
sideline & a football league's results, confirmations and table & Python, SQLite & 3 \\
skylark & a nature group's bird-sightings register with a paged feed & Node.js, JSON files & 3 \\
tally & a club election by ranked ballot & Node.js, JSON files & 3 \\
\bottomrule
\end{tabularx}
\end{table}

\section{The external benchmark}
\label{sec:external}

To compare with published RL web testers on their own ground, the final evaluation uses the
six subject applications of WebExplor~\cite{zheng2021webexplor}, which WebRLED's first
benchmark~\cite{gu2025webrled} also uses:
\begin{itemize}[noitemsep]
  \item Dimeshift, an expense tracker;
  \item Pagekit, a content management system, tested in its admin area;
  \item Splittypie, which splits bills;
  \item Phoenix, a Trello clone;
  \item Retroboard, a retrospective board;
  \item Petclinic, the Spring demonstration application.
\end{itemize}
They are open-source applications not written by this project. They run from the Docker
images of the DIG replication material~\cite{biagiola2019diversity}, which instrument the
client JavaScript with Istanbul, so branch coverage is measured the way the papers measured
it. To keep the comparison honest, nobody looked at these applications before the re-tuned
tester was frozen and the comparison pre-registered (Section~\ref{sec:extresults}).

\section{Admission}
\label{sec:admission}

Before any agent runs on an application, a harness must admit it mechanically. For each
fault, a scripted \emph{witness} trace must trigger the fault's assertion on the faulty build
and not on that fault's control build. The witness must replay deterministically and
minimise to a short reproduction. Benign traces must trigger nothing. A random walk on the
clean build must produce no candidate. The faulty build must differ from the clean build
\emph{only} by the declared patches; a diff audit checks this. The development suite adds a
data-model contract (seeded record counts per entity), a breadth contract (views, entities,
interaction patterns), sign-in checks, crash checks and a Docker build. An application that
fails any check is not used. A diversity check also confirms that the suite's applications
do not share code (the largest overlap was 3.2\%).

\section{Data split and independence}
\label{sec:split}

\begin{itemize}
  \item \textbf{Development} (used for tuning and design decisions): the six benchmark
        applications, the generated family, Gitea (a real open-source application, as a
        non-admin user), and the eight development suite applications.
  \item \textbf{Held-out}: the eight validation applications, used once, by the
        pre-registered evaluation of Chapter~\ref{chap.7}.
  \item \textbf{External}: the WebExplor and WebRLED subject applications, used once, by
        the final comparison.
\end{itemize}

The split was fixed before any agent ran on the suites. Development results are tuning data.
They are never presented as evidence of generalisation, and the held-out result belongs to
the version of the tester frozen before it ran.

\section{Metrics and statistics}
\label{sec:stats}

\paragraph{Outcome.} For a run with budget $B$, $T$ is the number of interactions until the
first verified fault. Every browser action counts, including replays and returns. A run that
finds nothing is right-censored at $T = B$.

\paragraph{Per arm.} The \emph{discovery rate} is the share of runs that verified at least
one fault. The \emph{restricted mean} of $T$ (RMST) averages $T$
over runs with censored runs at $B$. \emph{Distinct faults} counts how many of the seeded
faults an arm verified in at least one run. Clean-build control runs give the false-positive
count. For the tuning round two descriptive metrics were added. \emph{Unique failures} counts
distinct browser-error signatures (console, page and failed-request errors). \emph{JavaScript
coverage} is the share of the application's client-side code executed, measured with V8's
precise coverage through the Chrome DevTools Protocol, attached before the first navigation.

\paragraph{Comparisons.} Applications differ far more than seeds do, so the
\emph{application} is the unit of inference. For two arms, the per-application difference
in restricted mean $T$ is computed on matched (application, seed) pairs; the median across
applications summarises it, with a 95\% bootstrap interval resampling
applications. The stratified log-rank test is
secondary. The pre-registered decision rule:
\begin{itemize}[noitemsep]
  \item \emph{supported} if the median difference is negative (faster) and its interval
        excludes zero;
  \item \emph{refuted} if it is positive and the interval excludes zero;
  \item otherwise \emph{inconclusive}.
\end{itemize}

\paragraph{Pre-registration.} For every pre-registered experiment, the hypotheses, arms,
applications, budget, seeds, analysis code and decision rule were committed to the
repository, and the commit was tagged before any agent ran on the evaluation
applications~\cite{klees2018evaluating}. Deviations, such as rerunning a job lost to an
infrastructure failure, follow rules written in advance and are reported.

\section{Threats to validity}
\label{sec:threats}

\begin{itemize}
  \item \textbf{Same developer.} The test applications and their seeded faults were built
        by the same project as the tester. Independent planning of the validation suite and
        the external benchmark reduce, but do not remove, this threat.
  \item \textbf{Seeded faults.} Seeded faults are an incomplete proxy for real
        faults, and the oracle sees only what was seeded.
  \item \textbf{Scale.} Eight applications per evaluation leave wide intervals at the
        application level. Small effects cannot be told apart from noise.
  \item \textbf{Fault difficulty.} If the held-out faults are shallow, every arm finds them
        and differences compress. Chapter~\ref{chap.7} reports that this happened.
  \item \textbf{Tuning leakage.} Decisions made on development data can overfit it; the
        held-out and external evaluations exist to measure that.
\end{itemize}
